\documentclass[11pt]{article}
\usepackage[T1]{fontenc}
\usepackage{lmodern}
\usepackage[margin=1in]{geometry}
\usepackage{amsmath,amssymb,amsthm,mathtools}
\usepackage{microtype}
\usepackage[hidelinks]{hyperref}
\hypersetup{pdftitle={An n=8 counterexample to Karpenko's gamma-filtration conjecture},pdfauthor={Alec Kriebel}}
\newtheorem{theorem}{Theorem}
\newtheorem{lemma}[theorem]{Lemma}
\newtheorem{proposition}[theorem]{Proposition}
\newcommand{\CH}{\operatorname{CH}}
\newcommand{\OG}{\operatorname{OG}}
\newcommand{\IG}{\operatorname{IG}}
\newcommand{\SB}{\operatorname{SB}}
\newcommand{\ind}{\operatorname{ind}}
\newcommand{\gr}{\operatorname{gr}}
\newcommand{\C}{\mathbb C}
\newcommand{\Z}{\mathbb Z}
\newcommand{\Ftwo}{\mathbb F_2}
\title{An $n=8$ counterexample to Karpenko's\newline gamma-filtration conjecture}
\author{Alec Kriebel\\[2pt]\small Independent researcher\\
\small ORCID: \href{https://orcid.org/0009-0001-9320-500X}{0009-0001-9320-500X}}
\date{17 October 2026}
\begin{document}
\maketitle
\begin{abstract}
Karpenko's gamma-filtration nonmembership conjecture is false at $n=8$.
We spell out this consequence of his prior results: the later analytic
isotropy theorem for $\operatorname{Spin}(18)$, combined with the original
quadric-transfer argument and filtration theorem, gives
$e_8\in\gamma^5K_0(X_K)$ in the conjecture's precise maximal-Clifford-index
and Severi--Brauer function-field setup. An iterated Laurent-series field
provides an explicit characteristic-zero model. We also supply exact
arithmetic witnesses reproducing the previously published class valuation
underlying the earlier computational proof. This is a credited exposition
and verification of a literature-derived negative corollary; no firstness
claim is made for the corollary or its particular arithmetic witnesses.
\end{abstract}

\section{The original statement and the negative corollary}

Let $\varphi/F$ be a nondegenerate quadratic form of dimension $2n+1$,
with even Clifford algebra $A=C_0(\varphi)$ division of degree $2^n$.
Let $X=\OG_n(\varphi)$ parametrize maximal totally isotropic vector
subspaces, let $S=\SB(A)$ be the ordinary Severi--Brauer variety, and put
$K=F(S)$. We use integral $K_0$ throughout. For a smooth variety $V$,
$T^iK_0(V)$ is generated by coherent sheaves supported in codimension
at least $i$; $\gamma^iK_0(V)$ is the usual gamma filtration of its
augmented lambda ring.

Restriction identifies $K_0(X_K)$ with the split Grothendieck ring
\cite[\S4]{around}. Thus the structure-sheaf class of the split special
Schubert variety of codimension $n$ specifies a class
$e_n\in K_0(X_K)$. This definition does not assert that the corresponding
Schubert subvariety is defined over $K$. Karpenko's conjecture
\cite[Conjecture 3]{owr}, also \cite[Conjecture 4.7]{around}, asserts
\begin{equation}\label{eq:conjecture}
 e_n\notin\gamma^{\,n+1-\lfloor(n+1)/2\rfloor}K_0(X_K)
 \qquad(n\geq8).
\end{equation}

\begin{theorem}\label{thm:main}
Take
\[
 F=\C((t_1))\cdots((t_{16})),\qquad
 \varphi=\langle1,-t_1,\ldots,-t_{16}\rangle.
\]
Then $C_0(\varphi)$ is division of index $256$. With $X,S,K$ as above,
\[
                     e_8\in\gamma^5K_0(X_K).
\]
Consequently \eqref{eq:conjecture} fails at $n=8$.
\end{theorem}

The proof combines published results rather than asserting a new Spin
index theorem. The required generic fourth-Grassmannian index $8$ occurs
in author versions dated 2021 of \cite{oddspin}, is summarized in
\cite{indexes}, and is fully published in \cite{oddspin}. Karpenko's
later analytic argument \cite[Theorem 4.1 and Corollary 4.4]{isotropy}
provides the input used below without the earlier computer calculation.
The transfer from that input to the exact forbidden membership is already
given in \cite[following Conjecture 4.7]{around}. We make its application
and field identifications explicit. Priority for an earlier printed
statement of this combined conclusion has not been established; we do
not assert that the problem remained open until this note. Nor do we
determine exact gamma depth, an optimal replacement conjecture, or failure
at every $n\geq8$.

\section{The Clifford model and the isotropy input}

\begin{proof}[Verification of the model]
Give $F$ its composite rank-sixteen valuation, with value group
$\Gamma\simeq\Z^{16}$ and the sixteen values $v(t_i)$ a basis. If
$e_0,e_1,\ldots,e_{16}$ are Clifford generators for $\varphi$, then
$u_i=e_0e_i$ satisfy
\[
                  u_i^2=t_i,\qquad u_iu_j=-u_ju_i\quad(i\ne j).
\]
The ordered squarefree monomials $u_I$ form an $F$-basis of
$C_0(\varphi)$. Assign $u_I$ the value
$\frac12\sum_{i\in I}v(t_i)$ in $\frac12\Gamma$. Their values lie in
distinct cosets modulo $\Gamma$. Every nonzero linear combination
$\sum_I a_Iu_I$ therefore has a unique term of least value.
Multiplication of basis monomials respects these values and has nonzero
coefficient. The product of the two least terms is the unique least term
of a product, so two nonzero elements cannot have zero product.
The even Clifford algebra is central simple of dimension $2^{16}$;
absence of zero divisors makes it division, of degree $2^8$.
\end{proof}

Set
\[
 p=t_1\cdots t_{16},\qquad
 \psi=\varphi\perp\langle-p\rangle.
\]
The signed discriminant of $\psi$ is $p^2$, hence trivial. The standard
codimension-one Clifford and Grassmannian identifications
\cite[\S1]{around} give
\[
 C^+(\psi)\simeq C_0(\varphi)=A,
 \qquad X\simeq\OG^+_9(\psi),
\]
where either maximal component can be used. Since $K=F(\SB(A))$
splits $A$, the form $q=\psi_K$ has trivial discriminant and Clifford
invariant. No assertion that $q$ itself is hyperbolic is involved.

For a projective variety $W$ over a field $k$, its index is the positive
generator of the ideal of degrees of zero-cycles:
\[
 \ind(W)=\gcd\{[k(w):k]:w\text{ a closed point of }W\}.
\]

\begin{lemma}[Karpenko's isotropy bound]\label{lem:index}
If $q/k$ is an $18$-dimensional nondegenerate quadratic form in
characteristic zero with trivial discriminant and Clifford invariant,
then $\ind\OG_4(q)$ divides $8$.
\end{lemma}
\begin{proof}
The universal consequence of \cite[Theorem 4.1 and Corollary 4.4]{isotropy}
states that $q$ acquires Witt index at least $5$ over a finite extension
$L/k$ whose degree is not divisible by $16$. In particular
$\OG_4(q)(L)\ne\varnothing$, so its index divides $[L:k]$.
A diagonalization of $q$ and adjoining square roots gives a splitting
extension of $2$-power degree. Its Grassmannian therefore has $2$-power
index. These two divisibilities imply the assertion.
\end{proof}

The generic value in \cite{isotropy} is exactly $8$; only divisibility
by $8$ is used for our particular model. Alternatively the generic result
of \cite{oddspin} gives the same bound by generic-to-arbitrary divisibility
\cite[Proposition 2.1]{dkm}: the generic top-degree Chow image is contained
in the arbitrary one, so the arbitrary index divides the generic index.
This does not identify two different meanings of ``generic.''

\section{Transfer to the original gamma filtration}

\begin{proof}[Proof of Theorem \ref{thm:main}]
Write $Y$ for the projective quadric of $q=\psi_K$; thus $\dim Y=16$.
On a split quadric write $\ell_j$ for the structure-sheaf class of a
projective $j$-space, and put $h=1-[\mathcal O_Y(-1)]$.
The split-Clifford quadric $K$-theory identification used in
\cite[\S4]{around} gives these classes in $K_0(Y)$ and makes restriction
to a splitting field injective. They can be virtual classes: in
particular $\ell_0$ does not assert a $K$-rational point or membership
in $T^{16}K_0(Y)$.

Let $w$ be a closed point of $\OG_4(q)$, with residue field $L$ and degree
$d$. The corresponding isotropic vector $4$-space gives a projective
$3$-space on $Y_L$. Its class belongs to $T^{13}K_0(Y_L)$.
Finite pushforward preserves support codimension. After restriction
to a common splitting field it takes this class to $d\ell_3$.
Injectivity therefore gives
\[
                         d\ell_3\in T^{13}K_0(Y).
\]
By Lemma \ref{lem:index}, integral combinations of such degrees equal
$8$, and hence $8\ell_3\in T^{13}K_0(Y)$. An extension of degree
exactly $8$ is unnecessary.

The integral quadric identity
\cite[proof following Conjecture 4.7]{around} specializes to
\begin{equation}\label{eq:quadric}
          8\ell_3=\ell_0+h^{15}+2h^{14}+4h^{13}.
\end{equation}
For instance it follows by iterating
$2\ell_j=\ell_{j-1}+h^{16-j}$ for $j=1,2,3$.
The hyperplane terms lie in $T^{13}$, so \eqref{eq:quadric} gives
$\ell_0\in T^{13}K_0(Y)$, without division by $8$.

Use the even maximal component $X_K\simeq\OG^+_9(q)$ and its rank-nine
tautological bundle $U$. The incidence variety of an isotropic line
inside $U$ has projections
\[
             Y\xleftarrow{\ g\ }Z=\mathbf P(U)
                     \xrightarrow{\ f\ }X_K.
\]
Here $f$ is a projective bundle of relative dimension $8$, and $g$ is
smooth. The dimensions are $\dim X_K=36$, $\dim Z=44$, and $\dim Y=16$;
the relative dimension of $g$ is $28$. Flat pullback by $g$ preserves
support codimension, and proper pushforward by $f$ lowers it by at most
$8$. The split Schubert identity, valid in $K$-theory and transported
by restriction, is $e_8=f_*g^*(\ell_0)$
\cite[\S4, the $K$-theoretic version of (1.2)]{around}. Consequently
\[
                       e_8\in T^{13-8}K_0(X_K)=T^5K_0(X_K).
\]
Finally \cite[Theorem 4.3]{around} applies to our original odd form:
$A$ is division of maximal degree over $F$, and $K$ is precisely the
function field of its ordinary Severi--Brauer variety. It gives
$T^iK_0(X_K)=\gamma^iK_0(X_K)$ over this $K$ itself. Thus
$e_8\in\gamma^5K_0(X_K)$, whereas the term forbidden by
\eqref{eq:conjecture} at $n=8$ is $8+1-\lfloor9/2\rfloor=5$.
\end{proof}

The equality of filtrations in the last step is essential. Transporting
the underlying $K$-ring to a splitting field transports the gamma
filtration but need not make the topological filtration strict. In
particular, the split geometric codimension of the special Schubert
variety does not by itself settle this membership question.

\section{An exact certificate for the earlier finite calculation}

This section independently reproduces a published fact from
\cite[Appendix A]{oddspin}; it is additional verification, not a
dependency of the analytic proof above. Let
$V=\IG(5,16)=\operatorname{Sp}_{16}/P_5$, of dimension $45$,
and let $T$ be its rank-five tautological bundle. Write $x_1,\ldots,x_5$
for formal roots of $T^*$, $\pm y_1,\pm y_2,\pm y_3$ for those of
$Q=T^\perp/T$, and $u=\sum_i x_i$. The class in question is
\begin{equation}\label{eq:a}
 a=\prod_{\epsilon\in\{\pm1\}^3}
             (u+\epsilon_1y_1+\epsilon_2y_2+\epsilon_3y_3)
                 =2^8\widetilde z.
\end{equation}
Here $\widetilde z$ denotes the source's orbit product, interpreted in
the rational Chow ring for this identity. There are eight doubled orbit
factors. The prose in the inspected
author and journal versions of Appendix A prints $2^3\widetilde z$;
the product and \S\S4--5 require the normalization $2^8$ in
\eqref{eq:a}. We record this discrepancy without asserting an official
erratum or first identification. Simultaneously negating all roots
leaves this degree-eight product unchanged, so using $T^*$ is consistent.

Put $e_i=c_i(T^*)$, and let $h_j(x)$ be the complete homogeneous
symmetric polynomial. The exact sequences for $T\subset T^\perp$
give $c(T)c(T^*)c(Q)=1$, hence
\begin{equation}\label{eq:rank}
                       h_j(x_1^2,\ldots,x_5^2)=0\quad(j\geq4).
\end{equation}
Indeed these are the Chern classes $c_{2j}(Q)$ of a rank-six bundle.
If $q_i=e_i(y_1^2,y_2^2,y_3^2)$, expansion of \eqref{eq:a} is
\[
\begin{aligned}
 a={}&u^8-4q_1u^6+(6q_1^2-8q_2)u^4\\
     &+(-4q_1^3+16q_1q_2-64q_3)u^2+(q_1^2-4q_2)^2.
\end{aligned}
\]
Substitute $q_i=(-1)^ih_i(x^2)$ and call the resulting polynomial $A$.
Direct integral polynomial expansion gives
$A=128b+64h_4(x^2)$, where
\[
\begin{aligned}
b={}&-e_4^2-e_3e_5+2e_2e_3^2+5e_2^2e_4-2e_2^4
      +3e_1e_3e_4+4e_1e_2e_5\\
 &-9e_1e_2^2e_3-2e_1^2e_3^2-6e_1^2e_2e_4+7e_1^2e_2^3
      -e_1^3e_5\\
 &+7e_1^3e_2e_3+e_1^4e_4-5e_1^4e_2^2-e_1^5e_3+e_1^6e_2.
\end{aligned}
\]
Thus $a=128b$ in $\CH(V)$. A second polynomial identity is
\[
 b\equiv(e_1e_3+e_1^2e_2)h_4(x)+e_3h_5(x)+e_2h_6(x)
                      +e_1h_7(x)+h_8(x)\pmod2.
\]
Squaring modulo $2$ gives $h_j(x)^2=h_j(x^2)$, and
\eqref{eq:rank} gives $b^2=0$ in $\CH(V)/2$. Therefore
$a^3=2^{21}b^3$ is divisible by $2^{22}$ integrally.

The complementary degree is the exact integer
\begin{equation}\label{eq:pairing}
 \int_V a^3c_2(T^*)^9c_3(T^*)
 =14668659594398531584=2^{22}\cdot3497280977821.
\end{equation}
Here the complement has codimension $21=9\cdot2+3$, as required.
For an integer-only verification, $V$ is the zero locus of the regular
section of $\Lambda^2T^*$ induced by a fixed nondegenerate alternating form on
$\operatorname{Gr}(5,16)$, with class
$\prod_{i<j}(x_i+x_j)$. The usual Grassmannian coefficient formula
\cite{gysin} therefore evaluates the normalized degree as
\[
 [x_1^{15}x_2^{14}x_3^{13}x_4^{12}x_5^{11}]
 \left(b^3e_2^9e_3\prod_{i<j}(x_i+x_j)(x_i-x_j)\right)
                =6994561955642.
\]
Multiplication by $128^3$ gives \eqref{eq:pairing}. The archive evaluates
this by integral vertical-strip Pieri multiplication and, separately,
root-monomial coefficient extraction; exact fixed-point localization
provides further checks. These are established Schubert and Gysin
methods. The classes $c_2(T^*)$ and $c_3(T^*)$ here are dual-tautological
Chern classes, with no change to quotient-bundle one-row conventions.
Since $3497280977821$ is odd, the pairing prohibits
divisibility of $a^3$ by $2^{23}$. Together with the lower bound above,
this reproduces exact class valuation $\nu_2(a^3)=22$ in the free
integral Chow group. The valuation is prior work; priority for the
particular pairing and displayed identities remains unestablished.

\section{Verification and research disclosure}

The supporting verification archive contains self-contained exact
integer and finite-field checkers, source/version receipts, and
mathematical and priority review records. Independently implemented
routes reproduce both \eqref{eq:pairing} and the mod-$2$ coefficient
$[z^8x_7^4x_8^4]\prod_{I\subseteq\{5,6,7,8\}}
(z+\sum_{i\in I}x_i)=1$ used in \cite{isotropy}.
Recorded executions use ordinary Python and its optimization modes
\texttt{-O} and \texttt{-OO}. The original Maple program and its full
printed Schubert expansion were not replayed. These finite checks do
not reprove the cited Clifford, Chow, isotropy, filtration or Gysin
theorems. Third-party source bodies are omitted from the archive;
the inspected final author versions and bibliographic records are
distinguished in its source ledgers.

AI tools were used extensively in deriving and checking the argument,
literature retrieval and priority assessment, exact computation,
drafting, and internal adversarial review. This preprint is unrefereed;
it has not undergone external human peer review or proof-assistant
certification. The mathematical claims and their source dependencies
are presented for independent verification.

\begin{thebibliography}{9}
\bibitem{gysin}
L. Darondeau and P. Pragacz,
\emph{Universal Gysin formulas for flag bundles},
Internat. J. Math. \textbf{28}(11) (2017), 1750077.
\href{https://doi.org/10.1142/S0129167X1750077X}{doi:10.1142/S0129167X1750077X}.

\bibitem{dkm}
R. A. Devyatov, N. A. Karpenko and A. S. Merkurjev,
\emph{Maximal indexes of flag varieties for spin groups},
Forum Math. Sigma \textbf{9} (2021), e34.
\href{https://doi.org/10.1017/fms.2021.30}{doi:10.1017/fms.2021.30}.

\bibitem{owr}
N. A. Karpenko, \emph{Grothendieck ring of maximal orthogonal Grassmannian},
in \emph{Algebraic Cobordism and Projective Homogeneous Varieties},
Oberwolfach Report 5/2016, 226--227.
\href{https://doi.org/10.4171/OWR/2016/5}{doi:10.4171/OWR/2016/5}.

\bibitem{around}
N. A. Karpenko, \emph{Around 16-dimensional quadratic forms in $I_q^3$},
Math. Z. \textbf{285} (2017), 433--444.
\href{https://doi.org/10.1007/s00209-016-1714-x}{doi:10.1007/s00209-016-1714-x}.
Final author version, revised 4 June 2016:
\url{https://sites.ualberta.ca/~karpenko/publ/a16i5.pdf}.

\bibitem{oddspin}
N. A. Karpenko, \emph{On generic flag varieties for odd spin groups},
Publ. Mat. \textbf{67}(2) (2023), 743--756.
\href{https://doi.org/10.5565/PUBLMAT6722310}{doi:10.5565/PUBLMAT6722310}.

\bibitem{isotropy}
N. A. Karpenko, \emph{On isotropy fields of $\operatorname{Spin}(18)$-torsors},
St. Petersburg Math. J. \textbf{36}(1) (2025), 115--120;
Russian original: Algebra i Analiz \textbf{36}(1) (2024), 162--169.
\href{https://doi.org/10.1090/spmj/1848}{doi:10.1090/spmj/1848}.
Inspected author version dated 28 August 2023:
\url{https://sites.ualberta.ca/~karpenko/publ/spind18-04.pdf}.

\bibitem{indexes}
N. A. Karpenko and A. S. Merkurjev,
\emph{Indexes of generic Grassmannians for spin groups},
Proc. Lond. Math. Soc. (3) \textbf{125} (2022), 825--840.
\href{https://doi.org/10.1112/plms.12471}{doi:10.1112/plms.12471}.
\end{thebibliography}
\end{document}
